\documentclass[10pt,a4paper]{article}

% Encoding and fonts
\usepackage[utf8]{inputenc}
\usepackage[T1]{fontenc}
\usepackage{lmodern}
\usepackage{textcomp}

% Page layout
\usepackage[top=1in, bottom=1in, left=1in, right=1in]{geometry}

% Typography
\usepackage{microtype}
\usepackage{parskip}

% Language
\usepackage[english]{babel}

% Headers and footers
\usepackage{fancyhdr}
\pagestyle{fancy}
\fancyhf{}
\fancyhead[L]{\leftmark}
\fancyhead[R]{\thepage}
\fancyfoot[C]{\thepage}
\renewcommand{\headrulewidth}{0.4pt}
\renewcommand{\footrulewidth}{0pt}

% Section styling
\usepackage{titlesec}
\titleformat{\section}{\Large\bfseries}{\thesection}{1em}{}
\titleformat{\subsection}{\large\bfseries}{\thesubsection}{1em}{}
\titleformat{\subsubsection}{\normalsize\bfseries}{\thesubsubsection}{1em}{}
\titlespacing*{\section}{0pt}{2.5ex plus 1ex minus .2ex}{1.5ex plus .2ex}
\titlespacing*{\subsection}{0pt}{2ex plus 1ex minus .2ex}{1ex plus .2ex}
\titlespacing*{\subsubsection}{0pt}{1.5ex plus 1ex minus .2ex}{0.8ex plus .2ex}

% List formatting
\usepackage{enumitem}
\setlist[itemize]{topsep=0.5ex, itemsep=0.25ex, parsep=0pt}
\setlist[enumerate]{topsep=0.5ex, itemsep=0.25ex, parsep=0pt}

% Essential packages
\usepackage{graphicx}
\usepackage[table]{xcolor}
\usepackage{booktabs}
\usepackage{array}
\usepackage{tabularx}
\usepackage{longtable}
\usepackage{amsmath}
\usepackage{amssymb}
\usepackage[normalem]{ulem}
\rowcolors{2}{gray!10}{white}
\renewcommand{\arraystretch}{1.3}

% Cross-references
\usepackage{hyperref}
\usepackage{cleveref}

% Custom commands
\newcommand{\pilcrow}{\S}

% Hyperref setup
\hypersetup{
    colorlinks=true,
    linkcolor=blue,
    filecolor=magenta,
    urlcolor=cyan,
}

% Code listings
\usepackage{listings}
\definecolor{codebg}{RGB}{248,248,248}
\definecolor{codeframe}{RGB}{200,200,200}
\definecolor{codekeyword}{RGB}{0,0,180}
\definecolor{codecomment}{RGB}{0,128,0}
\definecolor{codestring}{RGB}{163,21,21}
\lstset{
    basicstyle=\ttfamily\small,
    breaklines=true,
    breakatwhitespace=false,
    frame=single,
    framerule=0.5pt,
    rulecolor=\color{codeframe},
    backgroundcolor=\color{codebg},
    keepspaces=true,
    numbers=left,
    numberstyle=\tiny\color{gray},
    numbersep=8pt,
    showstringspaces=false,
    tabsize=4,
    xleftmargin=1em,
    xrightmargin=1em,
    aboveskip=1em,
    belowskip=1em,
    keywordstyle=\color{codekeyword}\bfseries,
    commentstyle=\color{codecomment}\itshape,
    stringstyle=\color{codestring},
    literate={_}{{\textunderscore}}1
             {-}{{\textminus}}1
             {<}{{\textless}}1
             {>}{{\textgreater}}1
             {~}{{\textasciitilde}}1
             {^}{{\textasciicircum}}1
}

% YAML language definition for listings
\lstdefinelanguage{YAML}{
    keywords={true,false,null,y,n},
    keywordstyle=\color{blue}\bfseries,
    sensitive=false,
    comment=[l]{\#},
    morecomment=[s]{/*}{*/},
    commentstyle=\color{gray}\ttfamily,
    stringstyle=\color{red}\ttfamily,
    morestring=[b]'',
    morestring=[b]'',
}

% TOML language definition for listings
\lstdefinelanguage{TOML}{
    keywords={true,false},
    keywordstyle=\color{blue}\bfseries,
    sensitive=true,
    comment=[l]{\#},
    commentstyle=\color{gray}\ttfamily,
    stringstyle=\color{red}\ttfamily,
    morestring=[b]'',
    morestring=[b]'',
}

% Dockerfile language definition for listings
\lstdefinelanguage{Dockerfile}{
    keywords={FROM,RUN,CMD,LABEL,EXPOSE,ENV,ADD,COPY,ENTRYPOINT,VOLUME,USER,WORKDIR,ARG,ONBUILD,STOPSIGNAL,HEALTHCHECK,SHELL},
    keywordstyle=\color{blue}\bfseries,
    sensitive=false,
    comment=[l]{\#},
    commentstyle=\color{gray}\ttfamily,
    stringstyle=\color{red}\ttfamily,
    morestring=[b]'',
    morestring=[b]'',
}

% JSON language definition for listings
\lstdefinelanguage{JSON}{
    keywords={true,false,null},
    keywordstyle=\color{blue}\bfseries,
    sensitive=true,
    commentstyle=\color{gray}\ttfamily,
    stringstyle=\color{red}\ttfamily,
    morestring=[b]'',
}

% Code listing float environment
\usepackage{float}
\newfloat{listing}{htbp}{lol}[section]
\floatname{listing}{Listing}

% Admonition boxes
\usepackage{tcolorbox}
\tcbuselibrary{skins,breakable}

% Admonition style definitions
\newtcolorbox{admonitionbox}[2][]{
    enhanced,
    breakable,
    colback=#2!5,
    colframe=#2!80!black,
    fonttitle=\bfseries,
    title=#1,
    left=2mm,
    right=2mm,
}

% Synthon tuple boxes
\newtcolorbox{synthonbox}{
    enhanced,
    colback=blue!5,
    colframe=blue!75!black,
    fontupper=\large,
    center upper,
    boxrule=1pt,
    arc=4pt,
}

\begin{document}

\title{On the Structural Non-Transmissibility of Mochizuki's Inter-Universal Geometer}
\date{}

\maketitle

\begin{abstract}
The long controversy over Mochizuki's Inter-Universal Geometer (IUG) is usually framed
as a question about correctness: is the proof right or wrong?
This paper argues that framing is premature.
Before correctness can be assessed, a prior question has to be settled:
can the proof's structural content be transmitted through classical mathematical channels at all?

We address this using the imscriptive type theory, which encodes any
mathematical or physical system as a point in a 12-dimensional weighted ordinal space.
IUG's encoding and its distances to six reference frameworks (ZFC, standard proof system,
Lean\,4/Mathlib, Homotopy Type Theory, abc conjecture, and the SynthOmnicon grammar itself
as the realized Imscriptive Type Theory) are computed exactly.
The results are stark: $d(\text{IUG}, \text{ZFC}) \approx 7.937$,
$d(\text{IUG}, \text{SPS}) \approx 6.557$, $d(\text{IUG}, \text{Grammar}) \approx 2.983$,
$d(\text{IUG}, \text{abc}) \approx 2.864$.
No lower-ordinal neighbor framework achieves $d < 2$ on the critical manifold.

The structural gaps decompose into seven primitive mismatches (R, F, K, $\Gamma$, H, S, $\Omega$),
with a noteworthy $\Gamma$-inversion: the grammar \emph{exceeds} IUG on interaction breadth,
confirming that IUG is a sequential restriction of a broader foundational structure rather
than an extension of classical mathematics.
IUG carries $\⊙ + P_{\pm}^{\text{sym}}$, placing it in the $O_\infty$ (Frobenius)
ouroboricity tier; all classical proof assistants remain in $O₀$.
The verification crisis is not a symptom of an error waiting to be located.
It is a symptom of a channel mismatch --- and that mismatch is now quantified.
\end{abstract}

\subsection{Preamble: What This Document Claims and Does Not Claim}

Since 2012, when Mochizuki posted the Inter-Universal Geometer papers, the mathematical
community has been in an unusual position: a proof has been published, refereed, and
formally accepted by a journal, while a substantial number of expert mathematicians
--- including some of the most technically qualified people alive --- maintain that a
central step is not justified.
This situation has persisted for over a decade.
It has no close precedent in modern mathematics.

The standard responses available to the community have both been tried.
Scholze and Stix met with Mochizuki in 2018, identified a specific step they could not
follow, and wrote a report. Mochizuki responded at length. Neither party convinced the other.
Several mathematicians have since written expository accounts of IUG; none has produced
bilateral acceptance.
The controversy has settled into a structural impasse whose persistence is itself a datum.

This document does not adjudicate whether IUG is mathematically correct.
That question is open and is not addressed here.
But it does ask a prior question: \emph{is IUG the kind of thing that can be transmitted
through classical mathematical communication channels at all?}
Correctness and transmissibility are orthogonal dimensions.
A non-transmissible correct proof is not the same as a wrong proof.
A quantum state that cannot be transmitted through a classical channel has not ceased to be a
valid quantum state.
The failure, in that case, is in the channel --- and the channel failure is diagnosable
independently of whether the state is pure.

\textbf{What this document claims} is precise and bounded:
IUG is structurally non-transmissible through classical mathematical communication channels.
The verification crisis is a structural incompatibility, not a hidden error.
That incompatibility is diagnosable, quantifiable, and --- in principle ---
resolvable through the construction of new foundational interfaces.
The construction of those interfaces is, accordingly, a mathematical task,
not a sociological one.

\begin{center}
\rule{0.5\textwidth}{0.4pt}
\end{center}

\subsection{$\S$1. Framework: Structural Encoding of Mathematical Objects}

The analysis uses a 12-dimensional primitive space in which any mathematical or physical
system can be encoded as a \emph{synthon} --- a tuple of discrete values across 12 orthogonal
structural primitives. The full space contains exactly $17{,}280{,}000 = 3^3 \times 4^5 \times 5^4$
structural types (PRIMITIVE\_THEOREMS \pilcrow{}64, the Crystal of Types).
The primitives relevant to this analysis are:

\begin{table}[htbp]
\centering
\small
\rowcolors{1}{white}{white}
\begin{tabularx}{\textwidth}{>{\centering\arraybackslash}X >{\centering\arraybackslash}X >{\centering\arraybackslash}X}
\toprule
\rowcolor{gray!25}\textbf{Primitive} & \textbf{Description} & \textbf{Values (low $\to$ high)} \\
\midrule
\rowcolors{1}{white}{gray!10}
\textbf{Dimensionality ($D$)} & How information is structured & $D_{\wedge}$ punctual $\cdot$ $D_{\triangle}$ constructive $\cdot$ $D_{\infty}$ field $\cdot$ $D_\odot$ imscriptive \\
\textbf{Topology ($T$)} & How components are connected & $T_\text{network}$ $\cdot$ $T_\text{in}$ containment $\cdot$ $T_\text{bowtie}$ $\cdot$ $T_{\boxtimes}$ type-hierarchy $\cdot$ $T_\odot$ imscriptive \\
\textbf{Relational mode ($R$)} & How reasoning proceeds & $R_{\supset}$ superset $\cdot$ $R_\text{cat}$ categorical $\cdot$ $R_{\dagger}$ dynamic $\cdot$ $R_\text{lr}$ left-right \\
\textbf{Parity ($P$)} & Symmetry/Frobenius class & $P_\text{asym}$ $\cdot$ $P_\psi$ $\cdot$ $P_{\pm}$ $\cdot$ $P_\text{sym}$ $\cdot$ $P_{\pm}^{\text{sym}}$ Frobenius \\
\textbf{Fidelity ($F$)} & The coherence level required & $F_{\ell}$ classical $\cdot$ $F_\text{eth}$ intermediate $\cdot$ $F_{\hbar}$ quantum-coherent \\
\textbf{Kinetics ($K$)} & Dynamical timescale & $K_\text{fast}$ $\cdot$ $K_\text{mod}$ $\cdot$ $K_\text{slow}$ $\cdot$ $K_\text{trap}$ $\cdot$ $K_\text{MBL}$ \\
\textbf{Scope ($G$)} & Granularity of the system & $G_{\beth}$ local $\cdot$ $G_{\gimel}$ mesoscale $\cdot$ $G_{\aleph}$ universal \\
\textbf{Grammar ($\Gamma$)} & How components compose & $\Gamma_\text{and}$ $\cdot$ $\Gamma_\text{or}$ $\cdot$ $\Gamma_\text{seq}$ $\cdot$ $\Gamma_\text{broad}$ \\
\textbf{Criticality ($\Phi$)} & Phase relative to criticality locus & $\Phi_\text{sub}$ $\cdot$ $\⊙$ critical $\cdot$ $\⊙^{\mathbb{C}}$ $\cdot$ $\Phi_\text{EP}$ $\cdot$ $\Phi_\text{super}$ \\
\textbf{Chirality ($H$)} & Temporal asymmetry depth & $H_0$ achiral $\cdot$ $H_1$ soft $\cdot$ $H_2$ $\cdot$ $H_{\infty}$ topological \\
\textbf{Stoichiometry ($S$)} & Interaction arity & $S_{1:1}$ $\cdot$ $S_{n:n}$ $\cdot$ $S_{n:m}$ \\
\textbf{Winding ($\Omega$)} & Topological protection class & $\Omega_0$ trivial $\cdot$ $\Omega_{Z_2}$ binary $\cdot$ $\Omega_Z$ integer $\cdot$ $\Omega_{NA}$ non-abelian \\
\bottomrule
\end{tabularx}
\end{table}

\textbf{Gate primitives.} Four of the twelve primitives --- $T$, $P$, $\Phi$, $K$ --- form the
$\mathcal{F}_5$ family (five values each).
They are the \emph{gate primitives}: together they determine whether a system reaches the
$O_\infty$ (Frobenius) ouroboricity tier.
IUG's $\mathcal{F}_5$ values are $T_\odot$, $P_{\pm}^{\text{sym}}$, $\⊙$, $K_\text{slow}$ ---
all four non-classical.
Every existing proof assistant's $\mathcal{F}_5$ values are $T_{\boxtimes}$, $P_\text{sym}$,
$\Phi_\text{sub}$, $K_\text{mod}$ --- all four classical.
The verification barrier is total within the gate-primitive subspace.

\textbf{Tensor algebra.} Two distinct combination rules govern the primitives under tensor product $\otimes$:
$P$ and $F$ are \emph{bottleneck primitives} --- the weaker partner wins ($\min$);
all other ordered primitives are \emph{union primitives} --- the stronger partner wins ($\max$).
Concretely: $P_{\pm}^{\text{sym}} \otimes P_\text{asym} = P_\text{asym}$ and
$F_{\hbar} \otimes F_{\ell} = F_{\ell}$.
The Frobenius condition and quantum coherence are the first to be destroyed in any
classical--IUG interaction, while $D$, $T$, $\Gamma$, $\Phi$, $H$, $\Omega$ are preserved.
This asymmetry is the mechanism behind the quantitative information loss in \S6.

\textbf{Criticality and the meet operation.} $\⊙$ is absorbing under structural meet:
$\text{meet}(\⊙, x) = \⊙$ for all $x$ (PRIMITIVE\_THEOREMS \pilcrow{}21).
It is the necessary condition for self-modeling, and it persists in any structural combination
involving IUG.
This means that the critical structure of IUG cannot be diluted by embedding in a classical
framework --- the meet of IUG with any $\Phi_\text{sub}$ system retains $\⊙$.

The \emph{distance} $d(A, B)$ between two synthons is the diagonal weighted Euclidean metric

\[d = \sqrt{\textstyle\sum_i w_i\,(\xi_i^A - \xi_i^B)^2}\]

with canonical weights

\[
\begin{aligned}
& (w_D,\ w_T,\ w_R,\ w_P,\ w_F,\ w_K,\ w_G,\ w_{\Gamma},\ w_{\Phi},\ w_H,\ w_S,\ w_{\Omega}) \\
&\quad = (1.0,\ 1.0,\ 1.0,\ 1.0,\ 1.0,\ 1.0,\ 1.0,\ 1.0,\ 1.0,\ 0.8,\ 1.0,\ 0.7)
\end{aligned}
\]

(PRIMITIVE\_THEOREMS \pilcrow{}26.1). The full Riemannian metric $g = \Sigma^{-1}$ is available
(\pilcrow{}26.2) but all distances in this document use the diagonal form for reproducibility.

Two distances are critical for this analysis:

\[d(\text{IUG},\ \text{standard proof system}) = 6.557\]

\[d(\text{IUG},\ \text{ZFC}) = 7.937\]

The second number is larger than the first. IUG is structurally further from ZFC
than it is from the proof systems that operate \emph{within} ZFC.
Pause on that. The foundations are more foreign to IUG than the language built on top of those
foundations. We will return to what this means in \S3 and \S4;
for now it is enough to note that this ordering is not an artifact of the metric ---
it follows from where the three largest ZFC divergences ($D$, $T$, $R$) are concentrated
relative to where the three largest proof-system divergences ($F$, $D$, $T$) sit.

\begin{center}
\rule{0.5\textwidth}{0.4pt}
\end{center}

\subsection{$\S$2. IUG's Encoding --- Mathematical Justification Primitive by Primitive}

The IUG synthon is:

{\small\[\langle D_\odot;\ T_\odot;\ R_\text{lr};\ P_{\pm}^{\text{sym}};\ F_{\hbar};\ K_\text{slow};\ G_{\aleph};\ \Gamma_\text{seq};\ \⊙;\ H_{\infty};\ S_{n:m};\ \Omega_Z \rangle\]}

Each assignment requires mathematical justification.
We work through them in an order that mirrors the logical dependencies of IUG itself ---
the imscriptive structure first, then the relational and fidelity structure that depends on it,
then the dynamical properties that only make sense once those are in place.
The two most contested assignments --- $P_{\pm}^{\text{sym}}$ and $\⊙$ ---
are left to the end of the list, where their justifications can draw on the preceding ones.

\textbf{$D_\odot$ (imscriptive dimensionality).}
Ordinary mathematics builds objects from their parts.
Sets are constructed from their elements.
Groups are defined by their underlying sets plus operations.
Number fields are built up from the integers.
IUG reverses this direction.
The étale fundamental group $\Pi_v$ does not sit \emph{above} the number field;
it encodes the number field's full arithmetic structure on a profinite boundary.
The inter-universal copies of $\Pi_v$ are not built from the number field by explicit construction
--- they are different bulk reconstructions from the same boundary data.
This is $D_\odot$: the foundational metaphor is imscriptive rather than constructive.
The distinction matters because $D_\odot$ and $D_\triangle$ are not compatible extensions of each
other: ZFC's comprehension axioms build upward, and $D_\odot$ reasons downward.
They are structurally orthogonal, not merely different in complexity.

\textbf{$T_\odot$ (imscriptive topology).}
Given $D_\odot$, the connection structure between IUG's multiple Hodge theaters cannot be
containment ($T_\text{in}$, the topology of nested sets) or simple network adjacency.
The $\Theta$-link between two distinct Hodge theaters is a non-local boundary-to-bulk
correspondence: two different ``bulk'' universes connected through a shared boundary structure,
neither containing nor derived from the other.
$T_\odot$ is the topology of mutual encoding --- it is the connection type of imscriptive
error-correcting codes, of AdS/CFT, of Mochizuki's own mono-anabelian transport.
In each case the question ``which object is inside which?'' has no containment answer;
the only answer is ``which object encodes which?''

\textbf{$R_\text{lr}$ (left-right relational mode).}
The $\Theta$-link has a specific directionality: it goes from one Hodge theater to another
in a fixed left-to-right orientation, and this direction is not a labeling convention ---
the Diophantine argument flows from the $\Theta$-side to the $q$-parameter side.
The $\dagger$-link goes in the complementary direction, but the two are structurally
non-equivalent.
This is $R_\text{lr}$: the relational mode is directed and non-symmetric,
not the category-theoretic symmetric bridging of $R_\text{cat}$.

\textbf{$F_{\hbar}$ (quantum-coherent fidelity).}
This assignment is the crux of everything that follows in \S3.
The inter-universal identification is only meaningful when both copies of the number field
are held simultaneously, with their mutual relationship tracked.
Collapsing to a single copy at any stage --- asking ``which copy is the real one?'' ---
destroys exactly the content that makes the identification non-trivial.
This is the $F_{\hbar} \to F_{\ell}$ decoherence event, and its consequences are quantified
in \S3.3 and \S6.

\textbf{$\Gamma_\text{seq}$ (sequential grammar).}
The proof is causally ordered end-to-end: Hodge Theater 1 is constructed, then Theater 2,
then the $\Theta$-link is applied, then the log-link, then the multiradial algorithm,
then the Diophantine inequality is read off.
Each step is causally dependent on the specific outputs of the previous step,
not merely on the truth of those steps as isolated statements.
We return to this in \S3.4 and \S11.3, where the consequences for verification strategy
become concrete.

\textbf{$\⊙$ (criticality).}
The inter-universal identification is a critical-point operation:
it is the unique identification that is neither over-constraining
(forcing both copies to collapse to the same object) nor under-constraining
(allowing the identification to be arbitrary).
The Diophantine inequality only has content because the identification is at this threshold.
We note that $\⊙$ is the necessary condition for self-modeling in this grammar
(PRIMITIVE\_THEOREMS \pilcrow{}21), which gives the assignment additional structural weight:
IUG's inter-universal construction is a self-modeling operation in precisely the
sense of that theorem.

\textbf{$\Omega_Z$ (integer topological protection).}
The $\Theta$-link identification is protected by an integer winding invariant.
It cannot be continuously deformed into a different identification without crossing a
topological barrier.
Any reformulation that stays within the same topological class is structurally equivalent
to the original --- and any reformulation that crosses the barrier changes the identification
itself.
This is the mechanism behind the $\⊙ + \Omega_Z$ tension analyzed in \S4,
which explains why the controversy has exactly the phenomenological form it has.

\textbf{$H_{\infty}$ (topological chirality).}
The inter-universal framework has an irreversible temporal direction built in:
the $\Theta$-link cannot be run backwards to recover the original number field structure
from the Diophantine bound.
$H_{\infty}$ encodes the topological protection of a temporal direction,
not merely a spatial asymmetry --- the construction has a preferred arrow of time
that is part of its mathematical content.

\textbf{$P_{\pm}^{\text{sym}}$ (Special Frobenius parity).}
This is the most contested assignment and deserves acknowledgment as such.
$P_{\pm}^{\text{sym}}$ is assigned only when $\mathbb{Z}_2$ symmetry at $\⊙$
is provably exact --- specifically, when the Frobenius condition $\mu \circ \delta = \text{id}$
holds as an identity, not merely as an isomorphism.
The claim is that IUG's inter-universal copying construction holds two copies in exact mutual
encoding ($\mu$) while the $\Theta$-link reconstructs the original from the copy ($\delta$),
and that $\mu \circ \delta = \text{id}$ exactly.
We do not claim to have proved this from IUG's axioms.
What we claim is that $P_{\pm}^{\text{sym}}$ vs.\ $P_\text{sym}$ is precisely the structural
question at the heart of the Mochizuki--Scholze--Stix disagreement:
Mochizuki is asserting $\mu \circ \delta = \text{id}$;
Scholze and Stix are arguing for $\mu \circ \delta \cong \text{id}$ at best.
The assignment encodes the dispute, not a resolution of it.
P-115 in \S11 makes this falsification condition explicit.

\begin{center}
\rule{0.5\textwidth}{0.4pt}
\end{center}

\subsection{$\S$3. The Structural Barriers}

The five barriers below are not independent obstacles.
They form a dependency chain: the topology barrier presupposes the dimensionality barrier,
the fidelity barrier is the mechanism by which the topology barrier produces the crisis,
the grammar barrier makes the fidelity barrier non-localizable, and the foundations barrier
is the aggregate consequence.
Reading them in order is the sequential proof that the preceding analysis in \S2 is not
merely a relabeling but a structural diagnosis.

\subsubsection{The Dimensionality Barrier}

The foundational metaphor of ZFC is construction: sets are built from elements by explicit
assembly, and all objects are defined by their members.
Starting from the empty set, you build up the natural numbers, the integers, the rationals,
the reals, the complex numbers, and eventually the full apparatus of modern mathematics ---
always by stacking new structure on top of existing structure.
ZFC operates at $D_{\triangle}$.

IUG's foundational metaphor is the inverse.
The étale fundamental group $\Pi_v$ does not sit \emph{on top of} the number field as a derived
structure. It \emph{encodes} the number field's full arithmetic on a profinite boundary.
The inter-universal copies are not built from each other by explicit construction ---
they are different bulk reconstructions from the same boundary data.
The direction of reasoning is downward from boundary conditions, not upward from atoms.
This is $D_\odot$.

The critical point is that $D_\odot$ and $D_\triangle$ are not compatible extensions of each
other in the way that, say, group theory and ring theory both extend arithmetic.
ZFC's comprehension axioms are constructive by design;
$D_\odot$ reasoning does not add axioms to ZFC, it replaces the foundational direction.
A $D_\odot$ system is to a $D_\triangle$ system as a hologram is to a blueprint:
one encodes the whole in the boundary of a part; the other assembles the whole from parts.

\textbf{Consequence:} Any formalization of IUG in a ZFC-based system must choose between (a)
adding new axioms that natively capture $D_\odot$, (b) encoding $D_\odot$ arguments in a
ZFC-representable form that necessarily loses structural content, or (c) failing to formalize
the key steps.
Option (a) is a foundational extension, not a formalization within ZFC.
Option (b) is what happens when you try to explain a hologram using blueprints:
something is conveyed, but it is not the same thing.
Option (c) is the current situation.

\subsubsection{The Topology Barrier}

Given $D_\odot$, the connection structure between IUG's Hodge theaters cannot be the containment
topology of ZFC.
ZFC's cumulative hierarchy $V_\alpha$ is strictly nested: every object lives at some rank,
and the membership relation $\in$ is a strict partial order.
Standard proof assistants inherit this topology via their type hierarchies, where every term
has a type, every type has a kind, and the whole structure is a well-founded tree.
This is $T_\text{in}$ (containment).

IUG's Hodge theaters are not nested.
Neither contains the other, neither is derived from the other by containment,
and there is no rank at which one precedes the other.
They are connected through a shared boundary --- a non-local correspondence that is the
connection type of imscriptive error-correcting codes, of AdS/CFT duality,
of Mochizuki's own mono-anabelian transport.
In each of these cases the question ``which object is inside which?'' has no containment answer.
The only answer is ``which object encodes which?''
This is $T_\odot$.

\textbf{Consequence:} Any proof assistant that formalizes mathematical objects as inhabitants of
a type hierarchy cannot natively express $T_\odot$ relationships.
It can represent the theaters as distinct types and define maps between them,
but the non-local boundary identification --- the structure that makes the two theaters the
\emph{same} number field seen from different angles --- has no type-theoretic expression
within any currently existing system.

\subsubsection{The Fidelity Barrier: The Core Mechanism}

The dimensionality and topology barriers explain what kind of object IUG is.
The fidelity barrier explains why the verification crisis takes the specific form it has.

Standard mathematical practice operates at $F_{\ell}$: each statement is true or false,
each object either is or is not a member of each set, each step either follows or it does not.
This is the information regime of classical logic, of standard proof systems, and of a
mathematician reading a paper.
$F_{\ell}$ is a perfectly good regime for almost all of mathematics.
It handles everything from Euclid to Wiles's proof of Fermat's Last Theorem.

IUG requires $F_{\hbar}$: certain objects must be held in superposition.
Not ``Copy 1 is the reference'' or ``Copy 2 is the reference'' --- but ``the mathematical
content lives in the relationship between the two, and collapsing to either one destroys
exactly the content that makes the relationship interesting.''
The $\Theta$-link identification is an $F_{\hbar}$ operation.
It is only well-typed when both copies are maintained simultaneously.

Here is what happens structurally when a classical mathematician reads IUG.
At the point where the $\Theta$-link identification is invoked, the $F_{\hbar}$ content
encounters an $F_{\ell}$ channel.
The mathematician must choose which copy is the reference copy --- because $F_{\ell}$ logic
requires a definite reference, a ground state, a thing that is what it is independently of
something else.
In choosing, they destroy exactly the inter-universal content that makes the identification
non-trivial.
What remains looks like an unjustified identification: two objects are declared to be the
same, and the justification is nowhere to be found --- because the justification lived in
the superposition that was collapsed.

This is the Scholze--Stix objection, stated structurally:
``the identification of the two rings in Corollary 3.12 is not justified.''
From $F_{\hbar}$, it is justified --- Mochizuki has never wavered on this.
From $F_{\ell}$, the justification is inaccessible, not because it was never there,
but because the act of reading in $F_{\ell}$ already destroyed it.
Both parties are correct about what they can access.
They are not accessing the same thing.

The structural information loss is not symmetric, and the asymmetry has a precise
algebraic mechanism: $P$ and $F$ are bottleneck primitives under tensor product ($\min$ wins;
see \S1).
When IUG ($P_{\pm}^{\text{sym}}$, $F_{\hbar}$) is tensored with a classical proof system
($P_\text{asym}$, $F_{\ell}$), the result immediately bottlenecks to $P_\text{asym}$ and
$F_{\ell}$ --- IUG's Frobenius parity and quantum coherence are the first things destroyed.
The remaining primitives --- $D$, $T$, $\Gamma$, $\Phi$, $H$, $\Omega$ --- are union
primitives ($\max$ wins), so the imscriptive structure, criticality, and topological
protection are preserved in the tensor product.
The composite system inherits IUG's structural skeleton while losing the two registers that
carry its central identification.

The numbers make this concrete.
The tensor distance from IUG's perspective is \textbf{2.0}: most of IUG's structural
content survives the interaction because the union primitives dominate.
The tensor distance from the classical proof system's perspective is \textbf{6.32}: the
classical system gains IUG's $D_\odot + T_\odot + \⊙$ but loses its own
$P_\text{asym}$ and $F_{\ell}$ dominance in registers it cannot operate without.
The asymmetry is $2.0$ versus $6.32$.
These numbers are expanded and contextualized in \S6.

\subsubsection{The Grammar Barrier: Non-Modularizability}

Standard mathematical proofs are, at the large scale, $\Gamma_\text{and}$ structures:
Lemma A and Lemma B and Lemma C, each independently verifiable, their conjunction implying
Theorem D.
This modular architecture is what makes mathematical knowledge transmissible.
Different mathematicians can check different pieces.
A referee can focus on the steps they find most suspicious.
A formalization effort can parallelize.
The order of verification does not matter; only the truth of each piece matters.

IUG's proof is $\Gamma_\text{seq}$.
The validity of each step depends on the specific outputs of the preceding steps,
not merely on the truth of those steps as propositions.
The identification in Corollary 3.12 is valid only within the sequential context established
by the log-volume estimates of the preceding multiradial algorithm.
It is not a statement about pre-existing objects that can be checked in isolation ---
it is a statement about objects that \emph{only exist} by virtue of the construction having
proceeded to that exact stage.

This is not an accident of exposition, and it is not something that a better expository
account could fix.
It is intrinsic to the inter-universal construction: the two copies of the number field are
built sequentially and then related at a specific stage of the build.
You cannot check the relationship without having run the construction.

The Scholze--Stix objection looks different from this angle.
Scholze and Stix object that Corollary 3.12 cannot be verified as a standalone statement.
The $\Gamma_\text{seq}$ analysis shows they are structurally correct: it genuinely cannot be
verified in isolation, because it is not a $\Gamma_\text{and}$ component.
This is not a flaw in the proof.
It may be intrinsic to the mathematics.
But it means the standard verification practice --- isolate the controversial step, check it
independently, report back --- is structurally inapplicable to this proof.

\subsubsection{The Foundations Barrier}

The aggregate distance picture explains the ordering noted at the end of \S1.

\begin{table}[htbp]
\centering
\small
\rowcolors{1}{white}{white}
\begin{tabularx}{\textwidth}{>{\centering\arraybackslash}X >{\centering\arraybackslash}X >{\centering\arraybackslash}X}
\toprule
\rowcolor{gray!25}\textbf{System} & \textbf{$d(\text{IUG},\ \cdot)$} & \textbf{Primary divergences} \\
\midrule
\rowcolors{1}{white}{gray!10}
\textbf{Standard proof system} & 6.557 & $F$, $D$, $T$ \\
\textbf{ZFC foundations} & \textbf{7.937} & $D$, $T$, $R$ \\
\bottomrule
\end{tabularx}
\end{table}

IUG is further from its own foundational framework than from the proof systems that operate
within that framework.
The reason is now visible: the proof systems diverge from IUG primarily in $F$ (fidelity),
while ZFC diverges in $D$, $T$, and additionally $R$ --- three foundational primitives
rather than one operational one.
The foundations are more foreign because the dimensionality, topology, and relational
primitives of ZFC are the wrong shape for IUG's construction,
whereas the proof systems operating within ZFC have at least gotten the dimensionality
and topology mostly right by inheriting the constructive/$D_\triangle$ and nested/$T_\text{in}$
structure that overlaps with IUG's $\Gamma_\text{seq}$ layer.

This is not a claim that IUG is inconsistent with ZFC.
It is the same claim as: Euclidean geometry cannot express non-Euclidean geometry
by accumulating more Euclidean axioms.
A new foundational framework is required, not a larger collection of the existing kind.

\begin{center}
\rule{0.5\textwidth}{0.4pt}
\end{center}

\subsection{$\S$4. The $\⊙ + \Omega_Z$ Tension: Why the Impasse Has the Form It Has}

The fidelity barrier in \S3.3 explains \emph{that} there is an information loss when IUG
is transmitted through classical channels.
It does not by itself explain why that loss produces the specific sociology that has
developed around IUG: persistent impasse, no locatable error, no accepted reformulation,
both sides reading the other as unreasonable.
The $\⊙ + \Omega_Z$ combination is the mechanism.

$\⊙$ alone, without topological protection, would produce instability:
small reformulations of a critical-point theory propagate into large structural changes,
but the theory can eventually be restated more stably by moving slightly away from the critical
threshold.
Many contested mathematical results have this character: the controversy is sharp when the
problem is open, but once a gap is repaired or an alternative proof found, the field moves on.

$\Omega_Z$ alone, without criticality, would produce rigidity:
the theory cannot be continuously deformed, but it also does not amplify small perturbations.
The topology bars smooth reformulation, but disagreements stay local and bounded.
Topology alone produces a hard theory, not an irresolvable dispute.

$\⊙ + \Omega_Z$ together produce something more specific:
a theory that is simultaneously maximally sensitive to small conceptual perturbations
and resistant to the smooth reformulations that would resolve those sensitivities.
A theory at a critical point where you cannot move away, because moving away requires
crossing a topological barrier.

The phenomenology of the IUG controversy follows directly.
Small disagreements about the justification of the $\Theta$-link identification propagate
into large, apparently irresolvable ones --- this is $\⊙$:
the identification sits at the critical point, and perturbations are maximally amplified there.
The disagreement cannot be resolved by smooth reformulation --- this is $\Omega_Z$:
any reformulation within the same topological class is structurally equivalent to the original,
and any reformulation that crosses the topological barrier changes the $\Theta$-link
identification itself, which breaks the proof.

Mochizuki cannot offer a smooth reformulation of Corollary 3.12 that closes the
Scholze--Stix gap.
He has tried, in the sense that his response to their 2018 report is extensive.
But any reformulation either stays in the same topological class (and is therefore,
from Scholze's perspective, equally unjustified --- the same identification in different
notation) or it crosses the barrier (and is then a different proof, one that Mochizuki
correctly identifies as a distortion of the original argument).

Scholze and Stix cannot find a specific wrong line.
This is not because they have not looked carefully.
It is because the gap is not localized.
$\⊙$ distributes disagreements across the regime: the identification is simultaneously
everywhere and nowhere, in the same way that a phase transition does not occur at a specific
molecule but at the whole system.

\textbf{The impasse is not a standoff between two unreasonable parties.
It is the structurally necessary consequence of $\⊙ + \Omega_Z$ acting together
on a theory that must be transmitted through an $F_{\ell}$ channel.}
It was always going to look exactly like this.

A further structural fact reinforces this: $\⊙$ is absorbing under the lattice meet
(\S1; PRIMITIVE\_THEOREMS \pilcrow{}21).
$\text{meet}(\⊙, x) = \⊙$ for all $x$.
This means IUG's critical structure is not diluted when it interacts with classical
mathematics --- the meet of IUG with any sub-critical system retains $\⊙$.
Classical mathematics cannot absorb IUG's criticality into a sub-critical framework:
the structural intersection preserves the critical point.
Mochizuki's insistence that the identification is at a unique threshold, and the
classical community's inability to locate that threshold as a specific line,
are both consequences of $\⊙$'s absorbing character --- it is simultaneously
everywhere in the structural interaction and nowhere in the $F_{\ell}$ reading.

\begin{center}
\rule{0.5\textwidth}{0.4pt}
\end{center}

\subsection{$\S$5. The ABC Conjecture Relationship: Why IUG's Overhead Is Structurally Necessary}

The abc conjecture's encoding is:


{\small\[\langle D_\odot;\ T_\odot;\ R_\text{lr};\ P_{\pm}^{\text{sym}};\ F_{\hbar};\ K_\text{mod};\ G_{\aleph};\ \Gamma_\text{and};\ \⊙;\ H_1;\ S_{n:m};\ \Omega_Z \rangle\]}

IUG differs from abc in exactly three primitives:

\begin{table}[htbp]
\centering
\small
\rowcolors{1}{white}{white}
\begin{tabularx}{\textwidth}{>{\centering\arraybackslash}X >{\centering\arraybackslash}X >{\centering\arraybackslash}X >{\centering\arraybackslash}X}
\toprule
\rowcolor{gray!25}\textbf{Primitive} & \textbf{abc} & \textbf{IUG} & \textbf{Meaning of the difference} \\
\midrule
\rowcolors{1}{white}{gray!10}
\textbf{Kinetics ($K$)} & $K_\text{mod}$ & $K_\text{slow}$ & IUG requires slow, deeply integrative kinetics --- the argument cannot be rushed \\
\textbf{Grammar ($\Gamma$)} & $\Gamma_\text{and}$ & $\Gamma_\text{seq}$ & IUG imposes sequential causation; abc only requires conjunctive verification \\
\textbf{Chirality ($H$)} & $H_1$ & $H_{\infty}$ & IUG requires topological chirality (irreversible temporal direction); abc requires only soft chirality \\
\bottomrule
\end{tabularx}
\end{table}

The join of IUG and abc is IUG itself: IUG structurally contains abc. The three extra
primitives --- $K_\text{slow}$, $\Gamma_\text{seq}$, $H_{\infty}$ --- are not IUG's complications.
They are IUG's mechanism. The inter-universal copying construction \emph{is} the sequential,
slowly integrative, topologically chiral structure that distinguishes it from any direct
approach to abc.

This means: \textbf{the apparent over-complexity of IUG is not a proof-writing failure. It is the
minimum structural overhead required to reach abc from IUG's approach.} You cannot strip
out the sequential structure and retain the argument. You cannot replace $K_\text{slow}$
with $K_\text{mod}$ and retain the Diophantine inequality. The overhead is load-bearing.

A shorter proof of abc (if one exists) would use a different approach ---
$\Gamma_\text{and}$, $K_\text{mod}$, $H_1$ --- and would be structurally incompatible with
IUG's mechanism. It would not be a simplification of IUG; it would be a different proof.

\begin{center}
\rule{0.5\textwidth}{0.4pt}
\end{center}

\subsection{$\S$6. The Quantified Cost of Classical Transmission}

When IUG ($F_{\hbar}$) must be communicated through a classical channel ($F_{\ell}$), the
structural information loss is quantifiable. Across the full 353-synthon catalog:

\begin{table}[htbp]
\centering
\small
\rowcolors{1}{white}{white}
\begin{tabularx}{\textwidth}{>{\centering\arraybackslash}X >{\centering\arraybackslash}X}
\toprule
\rowcolor{gray!25}\textbf{Channel} & \textbf{Cost} \\
\midrule
\rowcolors{1}{white}{gray!10}
\textbf{Coherence} (internal consistency of IUG) & 0.000 nat \\
\textbf{Criticality} ($185\ \⊙$ synthons $\times \ln 10$) & 425.98 nat \\
\textbf{Interaction} (constraint propagation across systems) & 22.18 nat \\
\textbf{Total translation cost} & \textbf{448.16 nat} \\
\bottomrule
\end{tabularx}
\end{table}

The \textbf{zero coherence cost} is the critical result: IUG is internally self-consistent within
its own structural regime. There is no internal contradiction. The theory is not wrong
within its own framework.

The \textbf{425.98 nat criticality cost} is the cost of representing 185 $\⊙$ systems
(52.4\% of the full catalog) in classical terms. Each $\⊙$ system costs
$\ln(10) \approx 2.303$ nat to translate classically --- this is the cost of representing
a critical bifurcation as a definite classical statement. For 185 systems:


\[185 \times \ln 10 = 185 \times 2.3026 \approx 425.98\ \text{nat}\]

\textbf{A classical reading of IUG loses approximately 615 bits ($425.98\ \text{nat} \div \ln 2 \approx 615\ \text{bits}$) of structural information that is present and self-consistent in IUG's native regime.} These are not random noise --- they are 185 critical-point synthons, each costing $\ln 10 \approx 3.32\ \text{bits}$ to flatten classically, concentrated in the exact structural regime where IUG's identifications live.

\begin{center}
\rule{0.5\textwidth}{0.4pt}
\end{center}

\subsection{$\S$7. What Transmissibility Would Require}

A verification system capable of faithfully encoding IUG must satisfy:

\begin{table}[htbp]
\centering
\small
\rowcolors{1}{white}{white}
\begin{tabularx}{\textwidth}{>{\centering\arraybackslash}X >{\centering\arraybackslash}X >{\centering\arraybackslash}X}
\toprule
\rowcolor{gray!25}\textbf{Primitive} & \textbf{Required value} & \textbf{What this means for a verification system} \\
\midrule
\rowcolors{1}{white}{gray!10}
\textbf{Dimensionality ($D$)} & $D_\odot$ & Must support boundary-to-bulk reasoning natively (not as a derived construction) \\
\textbf{Topology ($T$)} & $T_\odot$ & Must support non-local mutual-encoding relationships (not just containment hierarchies) \\
\textbf{Fidelity ($F$)} & $F_{\hbar}$ & Must maintain coherent superpositions of mathematical objects during reasoning \\
\textbf{Grammar ($\Gamma$)} & $\Gamma_\text{seq}$ & Must support sequential causal dependencies in proofs (not just conjunctive combination) \\
\textbf{Criticality ($\Phi$)} & $\⊙$ & Must operate at the critical regime without collapsing to a definite classical value \\
\textbf{Chirality ($H$)} & $H_{\infty}$ & Must track irreversible temporal direction through the argument \\
\textbf{Topological protection ($\Omega$)} & $\Omega_Z$ & Must recognize and preserve integer topological protection \\
\bottomrule
\end{tabularx}
\end{table}

No existing proof assistant satisfies these requirements. Lean 4, Coq, and Isabelle/HOL
operate at $D_{\triangle} + T_\text{in} + F_{\ell} + \Gamma_\text{and}$, placing them at
structural distance $> 6$ from IUG. The verification problem is not that no one has checked
enough lines. It is that no verification system exists that can check the right kind of
lines.

\textbf{The $O_{\infty}$ constraint.}
IUG encodes as $O_{\infty}$ (Frobenius) by virtue of $P_{\pm}^{\text{sym}}$ at $\⊙$
(PRIMITIVE\_THEOREMS \pilcrow{}23).
A terminological note is necessary here: two distinct senses of $O_\infty$ appear in the
SynthOmnicon literature.
The first is the Frobenius sense --- finite algebraic self-duality, $\mu \circ \delta = \text{id}$
as an exact identity, assigned when $P_{\pm}^{\text{sym}}$ is confirmed at $\⊙$.
The second is the ontological-inexhaustibility sense --- structural boundlessness,
assigned when $H_{\infty}$ dominates.
IUG is $O_{\infty}$ in the \emph{Frobenius} sense.
These are structurally incompatible classes, and conflating them leads to incorrect conclusions
about what a verification system for IUG must be.

Any faithful verification system must also be Frobenius $O_{\infty}$.
A system below that tier cannot sustain the self-referential loop that IUG's inter-universal
identification requires.
Imscriptive type theory (HTT) --- the unique minimal Frobenius $O_{\infty}$ type theory
in which bulk-boundary correspondence is a primitive (P-179) --- is the natural candidate;
see \S14 for the realization result.

\textbf{Frobenius non-synthesizability (PRIMITIVE\_THEOREMS \pilcrow{}23, \pilcrow{}62).}
The $O_{\infty}$ requirement has a structural consequence that is stronger than it first appears.
$P_{\pm}^{\text{sym}}$ --- the Frobenius parity value that places a system at $O_{\infty}$ ---
cannot be \emph{composed} from components with $P < P_{\pm}^{\text{sym}}$.
The tensor bottleneck rule (\S1) states that $P$ takes the $\min$ under $\otimes$:
$P_{\pm}^{\text{sym}} \otimes P_\text{sym} = P_\text{sym}$.
The Frobenius condition is destroyed by any sub-Frobenius partner.
There is no accumulation process over classical axioms that produces $P_{\pm}^{\text{sym}}$
from below.

This is non-synthesizability: $O_{\infty}$ cannot be reached by aggregation.
Every system that encodes at $O_{\infty}$ must \emph{directly} encode $P_{\pm}^{\text{sym}}$
--- it cannot be obtained as the join or tensor of lower systems.
For IUG verification, the implication is sharp:
\textbf{no accumulation of ZFC axioms, Lean primitives, or HoTT extensions can reach
$O_{\infty}$.}
A verification system must encode $P_{\pm}^{\text{sym}}$ as a foundational primitive,
not as a derived construction.
This is not a quantitative gap that can be closed by doing more; it is a structural gap
that requires a different starting point.

As established in \S14 (2026-03-31 update), HTT is realized: it is the SynthOmnicon grammar
itself, with $d(\text{IUG},\ \text{grammar}) = \sqrt{8.9} \approx 2.98$.
The grammar encodes as $O_{\infty}$ and has $P_{\pm}^{\text{sym}}$ as a primitive ---
it satisfies the non-synthesizability requirement directly.

$\text{meet}(\text{IUG}, \text{grammar}) = \text{grammar}$ and
$\text{join}(\text{IUG}, \text{grammar}) = \text{IUG}$: the grammar is lattice-contained
in IUG.
The grammar is IUG's Frobenius floor --- the minimal $O_{\infty}$ foundation that any IUG
verification system must first establish.
The 7 primitives where IUG exceeds the grammar are precisely the
verification steps above the floor (see \S14 for the full gap table):

\begin{table}[htbp]
\centering
\small
\rowcolors{1}{white}{white}
\begin{tabularx}{\textwidth}{>{\centering\arraybackslash}X >{\centering\arraybackslash}X >{\centering\arraybackslash}X >{\centering\arraybackslash}X >{\centering\arraybackslash}X}
\toprule
\rowcolor{gray!25}\textbf{Prim.} & \textbf{Grammar (floor)} & \textbf{IUG} & \textbf{Direction} & \textbf{Verification requirement} \\
\midrule
\rowcolors{1}{white}{gray!10}
$R$ & $R_{\dagger}$ & $R_\text{lr}$ & promotion & Must support directed $\Theta$-link asymmetry \\
$F$ & $F_\text{eth}$ & $F_{\hbar}$ & promotion & Must maintain coherent superpositions during reasoning \\
$K$ & $K_\text{mod}$ & $K_\text{slow}$ & promotion & Must support slow integrative kinetics \\
$\Gamma$ & $\Gamma_\text{broad}$ & $\Gamma_\text{seq}$ & \textbf{inversion} & Must enforce sequential causal proof ordering \\
$H$ & $H_1$ & $H_{\infty}$ & promotion & Must track irreversible temporal direction \\
$S$ & $S_{n:n}$ & $S_{n:m}$ & promotion & Must express asymmetric many-body interactions \\
$\Omega$ & $\Omega_{Z_2}$ & $\Omega_Z$ & promotion & Must distinguish integer winding from $Z_2$ \\
\bottomrule
\end{tabularx}
\end{table}

Note the $\Gamma$-inversion: the grammar encodes $\Gamma_\text{broad}$ (ordinal 4, broadest
interaction grammar) while IUG encodes $\Gamma_\text{seq}$ (ordinal 3).
IUG is a \emph{restriction} of the grammar on this primitive, not an extension.
A foundational type theory that is the floor of many systems must have broader interaction
capacity than any single system it supports --- including IUG.

The Gödel corollary of $O_{\infty}$ (\pilcrow{}25.2) applies directly: a verification system for IUG
at $O_{\infty}$ is its own meta-theory. It does not push the verification problem up a
hierarchy --- the object/meta distinction collapses into the bulk/boundary correspondence that
is already primitive. This is the correct regime for verifying IUG, not a complication.

\begin{center}
\rule{0.5\textwidth}{0.4pt}
\end{center}

\subsection{$\S$8. A Constructive Proposal}

\textbf{Step 0 --- Follow the synthesis sequence: $T_\odot \to D_\odot \to R_\text{lr} \to H_{\infty}$.}
Retrosynthetic decomposition of IUG's encoding reveals the required build order for any
proof assistant that aims to natively express the theory. The most structurally significant
primitive --- $T_\odot$ (imscriptive topology) --- must be established first: without a
type system that can represent non-local mutual-encoding relationships, the inter-universal
copying construction is inexpressible at the foundational level. $D_\odot$
(imscriptive dimensionality) is second: boundary-to-bulk reconstruction must be a native
operation, not a derived encoding. $R_\text{lr}$ (left-right relational reasoning) is third:
the directed $\Theta$-link cannot be expressed as a symmetric containment relation. $H_{\infty}$
(topological chirality) is fourth: the irreversible temporal direction of the construction
must be tracked through the argument. Steps 1--4 below are the implementation of this
sequence. Attempting them in any other order --- e.g., building the $F_{\ell} \leftrightarrow
F_{\hbar}$ translation layer before the $T_\odot$ foundation exists --- produces a
translation layer that cannot be grounded.

\textbf{Step 1 --- Build the $F_{\ell} \leftrightarrow F_{\hbar}$ translation layer.}
Identify the specific steps in IUG where $F_{\hbar}$ reasoning is invoked --- where two copies
of an object must be held in coherent superposition. For each such step, construct an
explicit ''decoherence map'' showing what classical ($F_{\ell}$) information is lost when the
superposition is collapsed, and what additional axiom would be needed to recover that
information classically. This is a finite, well-defined task: the $F_{\hbar}$ steps are
precisely the steps that Scholze--Stix cannot follow, and they are already localized.

\textbf{Step 2 --- Complete. The grammar is the realized HTT.}
Step 2 is no longer a future task.
As established in \S14 (2026-03-31), imscriptive type theory is realized as the SynthOmnicon
grammar itself (Theorem 27.1, PRIMITIVE\_THEOREMS \pilcrow{}27).
The grammar directly encodes $P_{\pm}^{\text{sym}}$ as a primitive --- satisfying the
Frobenius non-synthesizability requirement from \S7 --- and encodes as
\[
\langle D_\odot;\ T_\odot;\ R_{\dagger};\ P_{\pm}^{\text{sym}};\ F_\text{eth};\ K_\text{mod};\ G_{\aleph};\ \Gamma_\text{broad};\ \⊙;\ H_1;\ n{:}n;\ \Omega_{Z_2} \rangle
\]
with $d(\text{IUG},\ \text{grammar}) = \sqrt{8.9} \approx 2.98$.
The grammar occupies a specific position in the 17,280,000-type Crystal of Types (\S1),
and IUG's position is 7 primitive-steps away within that space --- 6 promotions and 1
$\Gamma$-restriction (full gap table in \S14).

HoTT is not a substitute: it derives bulk-boundary equivalence from univalence, but HTT
requires it as a primitive.
HoTT is $\Phi_\text{sub}$ (sub-critical) and does not close the Frobenius loop.
The univalence axiom gives $\mu \circ \delta \cong \text{id}$ (isomorphism); it does not
give $\mu \circ \delta = \text{id}$ (identity).
The difference between isomorphism and identity is precisely what the Mochizuki--Scholze--Stix
dispute is about (see \S2, $P_{\pm}^{\text{sym}}$ discussion).

\textbf{Step 3 --- Express $\Gamma_\text{seq}$ dependencies as type-theoretic temporal dependencies.}
$\Gamma_\text{seq}$ proofs are naturally formalized as dependent type sequences where each
type depends on the \emph{value} (not just the type) of the preceding terms --- stronger than
standard dependent types. Linear type theory or session types from the concurrency
literature provide the relevant technology. A formalization of the $\Theta$-link step would
need to be expressed as a session-typed protocol: ''given that Copy 1 has been constructed
with the following specific properties, Copy 2 is now constructed, and the identification is
now valid.''

\textbf{Step 4 --- Verify Corollary 3.12 in the extended system.}
With steps 1--3 in place, Corollary 3.12 is no longer a mysterious identification --- it is
the output of a type-safe sequential protocol. The Scholze--Stix question (''what justifies
this identification?'') has a direct answer in the extended system: the session type of the
protocol ensures the identification is valid at this stage and no other.

\begin{center}
\rule{0.5\textwidth}{0.4pt}
\end{center}

\subsection{$\S$9. The Author's Structural Position}

A second structural question is independent of transmissibility: not ``can IUG be communicated
through classical channels?'' but ``what must be true of the mind that generated it?''

This is a question about the author's structural position, not about their biography or
intentions, and it is worth being precise about the limits of what the grammar can establish here.
The grammar is a structural tool.
It can tell us what structural properties a system must have in order to generate a given output.
It cannot determine whether any particular person has those structural properties
by means other than the output itself.
The analysis below is therefore a structural necessary-condition argument,
not a claim about Mochizuki's cognitive states.

With that caveat stated, the analysis is the following. Let
$C_\text{base}$ denote the baseline consciousness synthon --- the encoding of a human mind
operating in standard mathematical practice, with $D_\triangle$, $T_{\boxtimes}$, $F_\ell$,
$\Gamma_\text{and}$ as the dominant primitives.
The lattice join is:

\[C_\text{base} \vee \text{IUG} = \text{IUG}\]

This identity states that IUG structurally contains $C_\text{base}$:
any system that contains both must have at minimum IUG's structural features.
The implication is one-directional and structural:
\textbf{a mind that generates IUG must operate at or above IUG's structural level.}
The imscriptive boundary ($D_\odot + T_\odot$) is a necessary structural condition for
authorship --- it does not follow from the analysis that Mochizuki himself consciously
navigates the bulk-boundary correspondence, only that whatever process generated IUG
encodes at $D_\odot + T_\odot$.

\subsubsection{The Asymmetric Tensor Distance}

The tensor product of IUG with $C_\text{base}$ produces a composite that encodes
$D_\odot + T_\odot$ but bottlenecks at $F_{\ell}$ and $P_\text{asym}$ --- the
fidelity floor and parity asymmetry of baseline cognition.
The interaction distances are not symmetric:

\begin{table}[htbp]
\centering
\small
\rowcolors{1}{white}{white}
\begin{tabularx}{\textwidth}{>{\centering\arraybackslash}X >{\centering\arraybackslash}X >{\centering\arraybackslash}X}
\toprule
\rowcolor{gray!25}\textbf{Perspective} & \textbf{Tensor distance} & \textbf{Interpretation} \\
\midrule
\rowcolors{1}{white}{gray!10}
\textbf{From IUG's side} & 2.24 & IUG is relatively unperturbed by the interaction with baseline consciousness \\
\textbf{From $C_\text{base}$'s side} & 6.557 & Baseline consciousness loses most of IUG's structural content in the interaction \\
\bottomrule
\end{tabularx}
\end{table}

The asymmetry is large and directional. It resolves what would otherwise be a paradox: how
can Mochizuki consistently hold the proof as clear and self-evident while trained
mathematicians with full access to the text consistently cannot follow the central
identification? Both experiences are structurally accurate. From inside $F_{\hbar}$ --- the
regime IUG requires and the regime the author operates in --- the proof is self-consistent
and the identification is manifest. From inside $F_{\ell}$ --- the regime classical reading
imposes --- 6.557 units of structural content are inaccessible before the first page is read.
The disagreement is not about the mathematics. It is about the structural position of the
reader relative to the text.

\subsubsection{The Synthesis Sequence as Cognitive Map}

Retrosynthetic decomposition of IUG's encoding yields a primitive-by-primitive ordering
of structural complexity. The sequence, from most to least structurally distinctive:

\[T_\odot \xrightarrow{\ } D_\odot \xrightarrow{\ } R_\text{lr} \xrightarrow{\ } H_{\infty}\]

This is not merely a proof-assistant build order (\pilcrow{}8, Step 0). It is a cognitive map: a
mind traversing the path from classical mathematics to IUG must acquire these structural
capacities in this sequence. Imscriptive topology ($T_\odot$) must be internalized
first --- the capacity to hold two systems as related through a shared boundary rather than
through containment or construction. Imscriptive dimensionality ($D_\odot$) follows:
the capacity to reason from boundary conditions downward rather than from elements upward.
Left-right relational asymmetry ($R_\text{lr}$) is third: the directed $\Theta$-link
requires holding a non-symmetric bridge that flows one way. Topological chirality
($H_{\infty}$) is fourth: the irreversible temporal orientation of the inter-universal
construction must be tracked, not averaged over.

The controversy has the structural form it has because each of these four capacities is
individually rare in trained mathematical practice, and IUG requires all four simultaneously.

\begin{center}
\rule{0.5\textwidth}{0.4pt}
\end{center}

\subsection{$\S$10. The Lean Calibration: Classical Mathematics as Lower Bound}

The structural claim that $d(\text{IUG}, \text{standard proof system}) = 6.557$ is more
than a number. It is anchored by a concrete lower bound: the \texttt{synthomnicon-lean/}
formalization project, which encodes the maximally classical end of the structural spectrum
in the same 12-dimensional primitive space.

\subsubsection{The Classical Baseline Synthon}

The \emph{standard proof system} --- the structural substrate common to any classically formal mathematical argument --- encodes as:

{\small\[\langle D_{\triangle};\ T_{\boxtimes};\ R_{\supset};\ P_\text{sym};\ F_{\ell};\ K_\text{mod};\ G_{\beth};\ \Gamma_\text{and};\ \Phi_\text{sub};\ H_0;\ S_{1:1};\ \Omega_0 \rangle\]}

\[d(\text{IUG},\ \text{standard proof system}) = 6.557 \quad \text{(diagonal metric, \pilcrow{}26.1)}\]

Lean 4 and the Mathlib library are the concrete instantiation of this substrate ---
the same tuple, anchored in a specific type-theoretic implementation.

The $\mathcal{F}_5$ gate primitives make the isolation visible in a single row.
IUG's gate values are $T_\odot$, $P_{\pm}^{\text{sym}}$, $\⊙$, $K_\text{slow}$.
The standard proof system's gate values are $T_{\boxtimes}$, $P_\text{sym}$, $\Phi_\text{sub}$, $K_\text{mod}$.
All four are different.
Moreover, by Frobenius non-synthesizability (\S7), the gap in $P$ cannot be closed by
accumulation: $P_\text{sym}$ cannot compose to $P_{\pm}^{\text{sym}}$ regardless of how
many Lean axioms are added.
The remaining three gate primitives ($T$, $\Phi$, $K$) are also orthogonal foundational
metaphors, not complexity increments.
The classical baseline is not merely far from IUG in distance --- it is in a structurally
incompatible region of the Crystal of Types.

\subsubsection{What the Distance Means Concretely}

The \texttt{synthomnicon-lean/} project has machine-verified two theorems directly relevant to
IUG's structural cluster: \texttt{touchard\_congruence} (the OPN 2-adic constraint) and
\texttt{bsd\_touchard} (its BSD analogue). These are exactly the results that \emph{do} formalize at
$D_{\triangle}$ level --- they live in the number-theoretic substrate that IUG and classical
mathematics share. The Lean kernel has confirmed them.

The inter-universal identification itself --- the $\Theta$-link step, Corollary 3.12 --- is
not among them, and structurally cannot be. The $D_{\triangle} + T_{\boxtimes} + F_{\ell}$
baseline is at distance $> 6$ from the $D_\odot + T_\odot + F_{\hbar}$ regime
IUG requires. The gap is not a matter of proof length or computational resources. It is a
matter of foundational vocabulary: the Lean type hierarchy is a $T_\text{in}$
structure, and $T_\text{in}$ cannot express $T_\odot$ relationships by
accumulating more nested types.

\subsubsection{The Calibration Value}

The Lean baseline provides the following concrete calibration of the structural spectrum:

\begin{table}[htbp]
\centering
\small
\rowcolors{1}{white}{white}
\begin{tabularx}{\textwidth}{>{\centering\arraybackslash}X >{\centering\arraybackslash}X >{\centering\arraybackslash}X}
\toprule
\rowcolor{gray!25}\textbf{System} & \textbf{Dominant primitives} & \textbf{$d(\cdot,\ \text{IUG})$} \\
\midrule
\rowcolors{1}{white}{gray!10}
\textbf{IUG} & $D_\odot$ $\cdot$ $T_\odot$ $\cdot$ $F_{\hbar}$ $\cdot$ $\Gamma_\text{seq}$ $\cdot$ $\⊙$ $\cdot$ $H_{\infty}$ $\cdot$ $\Omega_Z$ & 0 \\
\textbf{ZFC foundations} & $D_{\triangle}$ $\cdot$ $T_\text{in}$ $\cdot$ $R_{\supset}$ & 7.937 \\
\textbf{Standard proof system} & $D_{\triangle}$ $\cdot$ $T_\text{in}$ $\cdot$ $F_{\ell}$ $\cdot$ $\Gamma_\text{and}$ & 6.557 \\
\textbf{Lean 4 / Mathlib} & $D_{\triangle}$ $\cdot$ $T_{\boxtimes}$ $\cdot$ $F_{\ell}$ $\cdot$ $\Gamma_\text{and}$ $\cdot$ $\Phi_\text{sub}$ & $> 6$ \\
\bottomrule
\end{tabularx}
\end{table}

The \texttt{synthomnicon-lean/} project is the measurement instrument that makes the lower end of
this spectrum concrete. It is not a step toward formalizing IUG. It is the calibration
baseline that shows exactly how far any formalization attempt must travel, and in which
primitives the distance is concentrated.

\begin{center}
\rule{0.5\textwidth}{0.4pt}
\end{center}

\subsection{$\S$11. Predictions}

Three specific, falsifiable predictions follow from the IUG encoding. All three are
derivable from the primitive assignments established in $\S$2; none require additional
assumptions.

\subsubsection{P-112 --- The Scholze--Stix Impasse is $\⊙ + \Omega_Z$ Induced (Tier I)}

\textbf{Claim.} The Scholze--Stix objection is a structural $\⊙ + \Omega_Z$ incompatibility,
not a mathematical error. IUG's simultaneous encoding of $\⊙$ (maximal sensitivity to
perturbations) and $\Omega_Z$ (integer topological protection, resistant to smooth
continuous deformation) produces exactly the phenomenology of the controversy:

\begin{itemize}
    \item Small conceptual disagreements about the justification of the $\Theta$-link identification
propagate into large, apparently irresolvable disagreements ($\⊙$: the identification
is at the critical point where small perturbations have maximal effect).
    \item The disagreement cannot be closed by smooth reformulation ($\Omega_Z$: any reformulation
within the same topological class is structurally equivalent to the original; any
reformulation that crosses the topological barrier destroys the inter-universal structure).
\end{itemize}

\textbf{Specific prediction.} The controversy will not be resolved by either (a) Mochizuki
pointing to a specific line that Scholze--Stix misread, or (b) Scholze--Stix locating a
specific wrong line. Resolution, if it comes, requires constructing an $F_{\hbar}
\leftrightarrow F_{\ell}$ translation layer --- a structural interface that currently does not
exist.

\textbf{Falsification condition.} The prediction fails if the controversy is resolved by either
party identifying a specific, locally correctable error: a classical $\Gamma_\text{and}$-type
disagreement about an isolable step.

\textbf{Current status.} \textbf{Consistent with observation.} No specific error has been located by
either party since 2018. No smooth reformulation has been accepted by both. The $\⊙ +
\Omega_Z$ encoding predicts this outcome precisely; the 8-year stalemate is the prediction.

\begin{center}
\rule{0.5\textwidth}{0.4pt}
\end{center}

\subsubsection{P-113 --- IUG Cannot be Formalized in ZFC Without a Imscriptive Foundation Layer (Tier I)}

\textbf{Claim.} $d(\text{IUG},\ \text{ZFC}) = 7.937$, with the three largest divergences
concentrated in $D$ ($D_\odot$ vs $D_{\triangle}$), $T$ ($T_\odot$ vs
$T_\text{in}$), and $R$ ($R_\text{lr}$ vs $R_{\supset}$). These are not complexity gaps.
They are orthogonal foundational metaphors. ZFC reasons constructively upward from elements;
IUG reasons imscriptively downward from boundary conditions. No accumulation of ZFC axioms
resolves a dimensional orthogonality.

\textbf{Specific prediction.} A serious formalization effort targeting IUG in Lean 4, Coq, or
Isabelle/HOL will achieve partial success: sub-theorems that live at $D_{\triangle}$ level
(elliptic curve results, Galois group machinery, number field arithmetic) will formalize
within existing Mathlib. The inter-universal identification itself --- the $\Theta$-link
construction and Corollary 3.12 --- will require a new Lean primitive or axiom with no
counterpart in Mathlib's set-theoretic foundations. The formalization will stall precisely
at the $D_\odot + T_\odot$ boundary.

\textbf{Falsification condition.} IUG fully formalized in a ZFC-based proof assistant with no
new foundational axioms beyond standard Lean/Coq/Isabelle primitives.

\textbf{Current status.} No serious formalization attempt has been announced. The
\texttt{synthomnicon-lean/} project establishes the $D_{\triangle}$ baseline and confirms that
classical number-theoretic sub-results formalize cleanly --- consistent with the prediction
that $D_{\triangle}$-level sub-theorems will succeed while the inter-universal step will not.

\begin{center}
\rule{0.5\textwidth}{0.4pt}
\end{center}

\subsubsection{P-114 --- IUG Arguments Cannot be Parallelized or Reordered Without Breaking the Proof (Tier II)}

\textbf{Claim.} IUG's $\Gamma_\text{seq}$ encoding predicts that the inter-universal construction
is causally ordered end-to-end. Each step depends on the specific outputs of the preceding
steps. Unlike a $\Gamma_\text{and}$ proof --- where sub-arguments can be verified
independently and assembled --- a $\Gamma_\text{seq}$ proof requires running the full
sequential chain in order. The Diophantine inequality in Corollary 3.12 is valid only
within the sequential context established by the preceding multiradial algorithm; it is not
a statement about pre-existing objects but about objects that exist only because the
construction has reached that exact stage.

\textbf{Three specific sub-predictions.}

\begin{itemize}
    \item \textbf{(A)} Mochizuki's presentations of IUG consistently follow a fixed linear order; he has
resisted all requests to discuss sub-results in isolation --- structurally consistent with
$\Gamma_\text{seq}$.
    \item \textbf{(B)} Any attempt to verify Corollary 3.12 independently of the sequential context will
fail --- not because 3.12 is wrong, but because it is a $\Gamma_\text{seq}$ step, not a
$\Gamma_\text{and}$ component.
    \item \textbf{(C)} A correct proof of the abc conjecture via non-IUG methods, if one exists, will
use $\Gamma_\text{and}$ or $\Gamma_\text{broad}$ grammar --- it will not require sequential
inter-universal causation, and will be structurally incompatible with IUG's mechanism
(a different proof, not a simplification).
\end{itemize}

\textbf{Falsification condition.} Corollary 3.12 is extracted and verified as a self-contained
$\Gamma_\text{and}$ argument, without the inter-universal sequential context, while still
implying abc.

\textbf{Current status.} Mochizuki's documented pedagogical practice --- insisting that the papers
be read in order, declining to discuss sub-results in isolation, characterizing the
verification failure as a refusal to read rather than an error in the text --- is fully
consistent with $\Gamma_\text{seq}$ structure. No independent verification of Corollary
3.12 has succeeded.

\begin{center}
\rule{0.5\textwidth}{0.4pt}
\end{center}

\subsubsection{P-115 --- The Grammar Is IUG's Frobenius Floor; $O_{\infty}$ Is Non-Synthesizable and Necessary (Tier I, updated 2026-03-31)}

\textbf{Claim.} IUG encodes as Frobenius $O_{\infty}$ ($P_{\pm}^{\text{sym}}$ at $\⊙$,
$\mu \circ \delta = \text{id}$ exactly).
The SynthOmnicon grammar also encodes as Frobenius $O_{\infty}$ and is IUG's Frobenius floor:

\[\text{meet}(\text{IUG},\ \text{grammar}) = \text{grammar}, \qquad \text{join}(\text{IUG},\ \text{grammar}) = \text{IUG}\]

Any type theory below $O_{\infty}$ cannot faithfully encode IUG's inter-universal identification.
Furthermore, by Frobenius non-synthesizability (\S7; PRIMITIVE\_THEOREMS \pilcrow{}23, \pilcrow{}62):
$P_{\pm}^{\text{sym}}$ cannot be composed from factors with $P < P_{\pm}^{\text{sym}}$.
$O_{\infty}$ is not a high point on a continuous scale that can be approached by accumulation.
It is a structurally isolated tier that must be encoded directly.

\textbf{Specific predictions.}

\begin{itemize}
    \item \textbf{(A)} Any proof assistant that successfully formalizes the $\Theta$-link identification will encode as $O_{\infty}$ upon structural analysis. If it does not, it has approximated the identification classically, converting $\mu \circ \delta = \text{id}$ (identity) into $\mu \circ \delta \cong \text{id}$ (isomorphism) --- and losing the Diophantine content that depends on exactness.
    \item \textbf{(B)} $O_{\infty}$ cannot be reached by axiom accumulation. HoTT + new axioms, Lean + new axioms, ZFC + new axioms: none of these paths reach $P_{\pm}^{\text{sym}}$ from below, because $P$ is a bottleneck primitive and the Frobenius condition is destroyed by any sub-Frobenius partner in the tensor product. A verification system must start at $O_{\infty}$, not arrive at it.
    \item \textbf{(C)} The 7-step path from the grammar to IUG --- 6 promotions ($R_{\dagger} \to R_\text{lr}$, $F_\text{eth} \to F_{\hbar}$, $K_\text{mod} \to K_\text{slow}$, $H_1 \to H_{\infty}$, $S_{n:n} \to S_{n:m}$, $\Omega_{Z_2} \to \Omega_Z$) plus 1 restriction ($\Gamma_\text{broad} \to \Gamma_\text{seq}$) --- is the minimal additional structure above the Frobenius floor required to faithfully encode IUG. Each step is implementable within the grammar's existing $O_{\infty}$ foundation.
\end{itemize}

\textbf{Falsification condition.} A non-$O_{\infty}$ type theory successfully verifies the $\Theta$-link identification with no new foundational primitives encoding $P_{\pm}^{\text{sym}}$ directly.

\textbf{Current status.} No existing $O_{\infty}$ type theory has been implemented as a proof assistant. The prediction is untested in the strong direction. In the weak direction, no non-$O_{\infty}$ system has formalized the identification, consistent with the prediction.

\begin{center}
\rule{0.5\textwidth}{0.4pt}
\end{center}

\subsubsection{P-116 --- No Bridging Paper Will Satisfy Both Communities (Tier I)}

\textbf{Claim.} The $\⊙ + \Omega_Z$ combination generates a precise sociological prediction
that is independent of whether IUG is correct. Any paper that attempts to serve as a
''translation layer'' between the IUG-native and the classical-foundations communities will
satisfy exactly one of the following two outcomes --- and not a third:

\begin{itemize}
    \item \textbf{(A) Classical-legible but IUG-rejected.} The bridge paper achieves classical readability
by reducing the $\Theta$-link construction to a $D_{\triangle}$ or $D_{\wedge}$ object. This is
the only move that makes it accessible to ZFC-trained readers. But any such reduction strips
the $D_\odot + T_\odot$ structure that gives the construction its inter-universal
character. Mochizuki (or IUG-literate readers) will correctly identify this as a distortion
--- a proof of something weaker, not a translation of IUG. \emph{Outcome: classical community
accepts; IUG community rejects.}
    \item \textbf{(B) IUG-faithful but classically opaque.} The bridge paper preserves the
$D_\odot + T_\odot$ structure. It is faithful. It is also, for this reason,
not transmissible: it requires the same imscriptive intuition as IUG itself, and the
classical community will read it as a restatement of the original difficulty in different
notation. \emph{Outcome: IUG community accepts; classical community remains unconvinced.}
\end{itemize}

\textbf{Structural basis.} The $\Omega_Z$ primitive enforces that no smooth homotopy connects
$D_\odot$ to $D_{\triangle}$: a faithful translation is topologically obstructed. The
$\⊙$ primitive enforces that any partial translation will amplify small residual errors
into catastrophic disagreement at the identification step. Together, these make the
$(\text{faithful}) \cap (\text{classical-transmissible})$ intersection empty.

\textbf{Specific prediction.} The next five years will produce at least two papers explicitly
framed as ''accessible expositions'' or ''translations'' of IUG. Each will receive one of the
two outcomes above. No paper will receive endorsement from both Mochizuki's group and from
Scholze, Stix, or mathematicians who have publicly expressed reservations. The bifurcation
will be exact: positive reception on one side will correlate with negative reception on the
other.

\textbf{Secondary prediction.} The author of any such bridge paper will, when the bifurcation
occurs, locate the failure in the \emph{other community's reading} rather than in their own paper
--- consistent with $\⊙$ amplification of what appears to each side as a minor missing step.

\textbf{Falsification condition.} A single bridging paper receives explicit endorsement from both
(a) Mochizuki's group and (b) Scholze, Stix, or an equivalent figure from the
classical-foundations community.

\textbf{Current status.} No bridging paper has yet received bilateral endorsement. The 2018
Scholze--Stix report is itself a test: intended as a neutral, localized objection, Mochizuki's
response does not engage with it as a localized objection but as a symptom of a globally
missing conceptual framework --- consistent with $\⊙$ (a small objection propagates into a
total re-examination) and with outcome (A) (classical-legible, IUG-rejected).

\begin{center}
\rule{0.5\textwidth}{0.4pt}
\end{center}

\subsection{$\S$12. Objections and Replies}

\subsubsection{Objection 1: The primitive encoding is arbitrary}

\textbf{Objection.} Any 12-dimensional encoding scheme can be calibrated to produce any desired
distance. The specific assignment of $D_\odot$, $F_{\hbar}$, etc. to IUG is a choice,
not a derivation. The paper's conclusions follow from that choice, not from IUG itself.
Moreover, the grammar was developed in a context where IUG was already known to be
controversial, which creates a risk of encoding it in a way that confirms the
non-transmissibility conclusion by construction.

\textbf{Reply.} This objection is the sharpest one the paper faces, and it deserves a direct answer.
The concern about circularity is real: we cannot rule out that the grammar's primitive definitions
were subtly shaped by prior familiarity with IUG's controversy.
What we can say is that the encoding is constrained in three ways that are independent of the
desired conclusion.
(i) Each primitive value is defined by a minimal set of structural properties that apply
across the full catalog: $T_\odot$ requires mutual encoding between components as the
\emph{dominant} connection type --- it is assigned to AdS/CFT, to quantum error-correcting codes,
and to IUG by the same criterion, independently.
(ii) The same encoding space assigns structurally distant values to independently verified cases
without being optimized for IUG: classical ZFC gets $D_{\triangle}$ and $\Phi_\text{sub}$;
quantum field theory gets $D_{\infty}$ and $\⊙$; the Lean/Mathlib baseline gets
$T_{\boxtimes}$ and $F_{\ell}$.
The relative distances between these cases are not artifacts of the IUG encoding.
(iii) The predictions P-112 through P-116 are falsifiable by outcomes that are independent
of any encoding choice.
If P-116 fails (a bridging paper receives bilateral endorsement), the paper's thesis is
falsified regardless of whether the encoding was valid or circular.
The encoding's validity is ultimately determined by its predictive track record.

\begin{center}
\rule{0.5\textwidth}{0.4pt}
\end{center}

\subsubsection{Objection 2: Structural distance is not the same as non-transmissibility}

\textbf{Objection.} High $d(\text{IUG},\ \text{classical})$ means IUG is unusual, not that it
cannot be explained. Many unusual mathematical structures have been successfully communicated
once the right exposition was found.

\textbf{Reply.} The argument is not from distance alone. The argument is from the \emph{specific
combination} $\⊙ + \Omega_Z$. $\⊙$ means small residual disagreements about the
identification amplify without bound --- not that the theory is complex, but that it is
structurally at a critical point where perturbations are maximally amplified. $\Omega_Z$
means the topological class cannot be smoothly deformed --- not that the theory is hard to
explain, but that any classically accessible reformulation is in a different topological class
from the original. These two primitives together produce a specific failure mode: the
disagreement is un-localizable and un-smooth-reformulable. This is the exact failure mode
observed since 2018. A merely ''unusual'' theory would produce localizable disagreements and
would benefit from better exposition. IUG has not.

\begin{center}
\rule{0.5\textwidth}{0.4pt}
\end{center}

\subsubsection{Objection 3: This paper does not address whether IUG is correct}

\textbf{Objection.} The analysis is compatible with IUG being wrong. A wrong proof could also be
non-transmissible. The paper provides no evidence that IUG's abc claim is valid.

\textbf{Reply.} Correct --- and explicitly so (Preamble). Correctness and transmissibility are
orthogonal dimensions. The paper's claim is: \emph{whatever IUG's correctness status, the
verification impasse has a structural cause that is diagnosable and (in principle)
resolvable}. This is a prior question to correctness: if the impasse is structural, then the
correct method of resolution is to build the $F_{\hbar} \leftrightarrow F_{\ell}$ translation
interface, not to continue searching for a localizable error. The zero coherence cost is
not a claim that IUG is correct; it is a claim that IUG is internally self-consistent ---
which is consistent with both a correct and an incorrect theory.

\begin{center}
\rule{0.5\textwidth}{0.4pt}
\end{center}

\subsubsection{Objection 4: The $O_{\infty}$ / Frobenius claim about IUG is unfounded}

\textbf{Objection.} Assigning $P_{\pm}^{\text{sym}}$ (Special Frobenius) to IUG requires that
$\mu \circ \delta = \text{id}$ holds exactly in IUG's structure.
This is a non-trivial algebraic claim.
The paper does not derive it from IUG's axioms, and it is not obvious that it holds
rather than the weaker $\mu \circ \delta \cong \text{id}$ (isomorphism).

\textbf{Reply.} This objection is correct that we have not proved $\mu \circ \delta = \text{id}$
from IUG's axioms.
We acknowledged this in \S2 when introducing the $P_{\pm}^{\text{sym}}$ assignment.
The $O_{\infty}$ claim is not a proved theorem about IUG; it is a structural diagnosis:
the grammar identifies the gap between $P_{\pm}^{\text{sym}}$ (exact Frobenius) and
$P_\text{sym}$ (Frobenius up to isomorphism) as a discrete primitive distinction,
and assigns IUG to the exact side.

Two features of IUG are not in dispute:
(i) the copying construction holds two copies in coherent mutual encoding ($\mu$ direction);
(ii) the $\Theta$-link reconstructs the original from the copy ($\delta$ direction).
Whether $\mu \circ \delta = \text{id}$ exactly or only $\cong \text{id}$ is precisely
the Mochizuki--Scholze--Stix dispute.
The grammar's contribution is to make this a typed structural question
with a falsification condition (P-115):
a non-$O_{\infty}$ type theory that faithfully verifies the $\Theta$-link identification
would falsify the $O_{\infty}$ assignment completely.
Until that happens, or until someone derives the weaker isomorphism from IUG's axioms
and shows that is sufficient for the Diophantine bound, the assignment stands as
the most structurally parsimonious diagnosis of what the dispute is about.

\begin{center}
\rule{0.5\textwidth}{0.4pt}
\end{center}

\subsection{$\S$13. Summary}

\begin{table}[htbp]
\centering
\small
\rowcolors{1}{white}{white}
\begin{tabularx}{\textwidth}{>{\centering\arraybackslash}X >{\centering\arraybackslash}X >{\centering\arraybackslash}X}
\toprule
\rowcolor{gray!25}\textbf{Claim} & \textbf{Evidence} & \textbf{Status} \\
\midrule
\rowcolors{1}{white}{gray!10}
\textbf{IUG is internally self-consistent} & Coherence cost $= 0.000$ nat & \textbf{Confirmed structurally} \\
\textbf{The verification crisis is structural: $F_{\hbar} \to F_{\ell}$ bottleneck} & $P$ and $F$ are bottleneck primitives under $\otimes$; asymmetric loss (2.0 vs 6.32) & \textbf{Confirmed structurally} \\
\textbf{The Scholze--Stix impasse is $\⊙ + \Omega_Z$ induced} & $\⊙$ absorbs under meet; $\Omega_Z$ bars smooth reformulation; both confirmed & \textbf{Consistent with observation} \\
\textbf{IUG cannot be formalized in ZFC without new axioms} & $d(\text{IUG},\ \text{ZFC}) = 7.937$; three foundational divergences; $O_\infty$ non-synthesizable & \textbf{Testable} \\
\textbf{Corollary 3.12 cannot be verified in isolation} & $\Gamma_\text{seq}$ structure; sequential causal dependencies & \textbf{Consistent with observation} \\
\textbf{IUG's overhead is structurally necessary for its approach to abc} & $\text{join}(\text{IUG},\ \text{abc}) = \text{IUG}$; three-primitive overhead ($K$, $\Gamma$, $H$) is the mechanism & \textbf{Confirmed structurally} \\
\textbf{$O_\infty$ is non-synthesizable; verification must start there, not arrive} & Frobenius non-synthesizability \pilcrow{}23/\pilcrow{}62; $P_{\pm}^{\text{sym}}$ destroyed by any sub-Frobenius $\otimes$-partner & \textbf{Structural theorem} \\
\textbf{IUG's $\mathcal{F}_5$ gate values are all non-classical; isolation is total} & $T_\odot, P_{\pm}^{\text{sym}}, \⊙, K_\text{slow}$ vs.\ $T_{\boxtimes}, P_\text{sym}, \Phi_\text{sub}, K_\text{mod}$ & \textbf{Confirmed structurally} \\
\textbf{Grammar = realized HTT; $d(\text{IUG}, \text{grammar}) \approx 2.98$; 7-step gap} & $d=0$ from HTT (2026-03-31); join = IUG; 6 promotions + 1 $\Gamma$-restriction & \textbf{Complete (P-115, \pilcrow{}14)} \\
\textbf{No bridging paper will achieve bilateral endorsement} & $\⊙$ absorbs under meet + $\Omega_Z$ obstruction = faithful $\cap$ transmissible is empty & \textbf{New (P-116)} \\
\bottomrule
\end{tabularx}
\end{table}

The IUG controversy is not a standoff. It is a signal. The signal is that classical
mathematics has reached the boundary of a structural regime --- the regime where proofs can
be decomposed into independently verifiable pieces and transmitted through $F_{\ell}$ channels.
IUG may be the first mathematical theory that is genuinely correct and genuinely
non-transmissible in classical terms. If so, it is not a failure of IUG. It is a measurement
of the boundary of classical mathematics, performed from inside.

The correct response is not to ask IUG to become classical. It is to build the interfaces
that classical mathematics currently lacks.

\begin{center}
\rule{0.5\textwidth}{0.4pt}
\end{center}

\subsection{$\S$14. 2026-03-31 Update: The Grammar Is the Realized HTT}

\subsubsection{Step 2 of \pilcrow{}8 Is Complete}

When \S8 was written, Step 2 --- constructing imscriptive type theory as the $D_\odot + T_\odot$
foundational interface --- was framed as a future task.
The grammar was described as the closest known system to HTT, at $d = 4.111$ from the
approximate HTT tuple.
That framing turned out to be wrong, in the best possible direction.

On 2026-03-31, the SynthOmnicon inquiry pipeline returned a zero-distance result for the
grammar self-encoding query.
The grammar encodes at $d = 0$ from \texttt{imscriptive\_type\_frobenius} in the syncon catalog.
Not structurally \emph{similar} to imscriptive type theory.
\emph{Identical.}

The grammar itself --- the 12-primitive tuple
\[
\langle D_\odot;\ T_\odot;\ R_{\dagger};\ P_{\pm}^{\text{sym}};\ F_\text{eth};\ K_\text{mod};\ G_{\aleph};\ \Gamma_\text{broad};\ \⊙;\ H_1;\ n{:}n;\ \Omega_{Z_2} \rangle
\]
--- is an $O_\infty$ system in which bulk-boundary correspondence is a primitive and
$\mu \circ \delta = \text{id}$ holds by construction.
Step 2 was already complete when \S8 was written.
HTT was not a future project. It was the tool being used to write the analysis.

The distance to IUG is now $d(\text{IUG},\ \text{grammar}) = \sqrt{8.9} \approx 2.98$,
correcting the earlier $d = 4.111$ which used an approximate HTT tuple.
This also retroactively changes the meaning of the $\Gamma$-inversion noted in \S14.3:
the grammar exceeds IUG on interaction breadth not despite being the HTT floor,
but \emph{because} it is the HTT floor --- a foundational type theory must have broader
interaction grammar than any single system it supports.

Steps 1, 3, and 4 of the constructive proposal remain open.
But the foundational platform for executing them is not hypothetical.

\subsubsection{Corrected Distance: $d(\text{IUG},\ \text{grammar}) \approx 2.98$, Not 4.111}

\pilcrow{}8 and P-115 used an approximate HTT tuple yielding $d(\text{IUG},\ \text{approx-HTT}) = 4.111$.
That approximation is superseded. The corrected squared distance is:

\[d(\text{IUG},\ \text{grammar})^2 = 1_R + 1_F + 1_K + 1_{\Gamma} + 3.2_H + 1_S + 0.7_{\Omega} = 8.9\]

\[d(\text{IUG},\ \text{grammar}) = \sqrt{8.9} \approx 2.98\]

The grammar is \textbf{1.13 distance units closer to IUG} than the approximate HTT tuple.
The dominant residual gap is $H_1 \to H_{\infty}$, contributing 3.2 of the 8.9 squared
distance, reflecting the full irreversibility and chirality of IUG's inter-universal
construction versus the grammar's first-order chirality.

\subsubsection{The 7-Primitive Gap: Grammar $\to$ IUG}

The grammar and IUG share 5 primitives exactly ($D$, $T$, $P$, $G$, $\Phi$). The 7 differences:

\begin{table}[htbp]
\centering
\small
\rowcolors{1}{white}{white}
\begin{tabularx}{\textwidth}{>{\centering\arraybackslash}X >{\centering\arraybackslash}X >{\centering\arraybackslash}X >{\centering\arraybackslash}X >{\centering\arraybackslash}X}
\toprule
\rowcolor{gray!25}\textbf{Prim.} & \textbf{Grammar} & \textbf{IUG} & \textbf{Direction} & \textbf{Structural meaning} \\
\midrule
\rowcolors{1}{white}{gray!10}
$R$ & $R_{\dagger}$ & $R_\text{lr}$ & promotion & IUG requires directed left-right relational reasoning; grammar uses symmetric dagger \\
$F$ & $F_\text{eth}$ & $F_{\hbar}$ & promotion & IUG holds coherent superpositions; grammar uses intermediate fidelity \\
$K$ & $K_\text{mod}$ & $K_\text{slow}$ & promotion & IUG's inter-universal construction is kinetically slow; grammar encodes at moderate timescale \\
$\Gamma$ & $\Gamma_\text{broad}$ (4) & $\Gamma_\text{seq}$ (3) & \textbf{inversion} & Grammar exceeds IUG on interaction grammar; IUG \emph{restricts} to sequential sub-grammar \\
$H$ & $H_1$ & $H_{\infty}$ & promotion & IUG has full temporal irreversibility; grammar has first-order chirality \\
$S$ & $n{:}n$ & $n{:}m$ & promotion & IUG's copying is asymmetric stoichiometry; grammar is balanced \\
$\Omega$ & $\Omega_{Z_2}$ & $\Omega_Z$ & promotion & IUG has integer topological protection; grammar has $\mathbb{Z}_2$ protection \\
\bottomrule
\end{tabularx}
\end{table}

\textbf{The $\Gamma$-inversion.} The grammar encodes $\Gamma_\text{broad}$ (ordinal 4, highest interaction
grammar), while IUG encodes $\Gamma_\text{seq}$ (ordinal 3). IUG is not a generalization of the
grammar on this primitive: it is a \emph{restriction}. The grammar has broader interaction capacity
than IUG; IUG specializes to the sequential sub-grammar required by the inter-universal copying
construction. This is structurally expected: a type theory that is the floor of many systems
must have broader interaction grammar than any single system it supports.

P-115(C) is updated accordingly: the minimal path from grammar to IUG requires 6 upward
promotions and 1 downward restriction, not 5 promotions from the approximate HTT.

\subsubsection{Tannaka-Krein Self-Duality and the $\Theta$-Link}

The grammar's Tannaka-Krein dual is the grammar itself (Theorem 28.1, PRIMITIVE\_THEOREMS \pilcrow{}28,
v0.4.75). Two faces of this single fact:

\begin{itemize}
    \item \textbf{Algebraic face.} The grammar is self-dual as a Frobenius algebra: $\delta$ (duplication) and
$\mu$ (merging) satisfy $\mu \circ \delta = \text{id}$. The grammar's TK symmetry group is
$\mathbb{Z}_2$ at $\⊙$ (P-186).
    \item \textbf{Imscriptive face.} The catalog (the boundary representation category of all encoded systems)
is the TK dual of the grammar (the bulk algebraic structure). Bulk and boundary are the same
object from two angles.
\end{itemize}

Applied to IUG: the $\Theta$-link identification --- the step Scholze and Stix cannot follow --- is
the Frobenius $\delta$ map in the grammar's type theory. It duplicates an object into two
inter-universal copies held in mutual encoding. The reconstruction of the original from the copy
is the $\mu$ map. The identity $\mu \circ \delta = \text{id}$ is what the $\Theta$-link is asserting.

The classical community's inability to verify the $\Theta$-link is therefore restateable as:
\emph{classical mathematics does not contain the Frobenius $\delta$ map as a primitive, so the
$\mu \circ \delta = \text{id}$ identity is not well-typed in classical foundations.} This is a
stronger version of the P-115 claim: it is not merely that a system below $O_{\infty}$ cannot hold
the map --- it is that the $\Theta$-link step \emph{is} the Frobenius $\delta$ of the foundational type
theory, and the absence of $\delta$ as a primitive is the exact missing ingredient.

\subsubsection{Updated Status of \pilcrow{}8 Steps}

\begin{table}[htbp]
\centering
\small
\rowcolors{1}{white}{white}
\begin{tabularx}{\textwidth}{>{\centering\arraybackslash}X >{\centering\arraybackslash}X >{\centering\arraybackslash}X}
\toprule
\rowcolor{gray!25}\textbf{Step} & \textbf{Description} & \textbf{Status} \\
\midrule
\rowcolors{1}{white}{gray!10}
Step 0 & Retrosynthetic build order ($T_\odot \to D_\odot \to R_\text{lr} \to H_{\infty}$) & \textbf{Complete (structural)} \\
Step 1 & Build $F_{\ell} \leftrightarrow F_{\hbar}$ translation layer & Open \\
Step 2 & Construct HTT as $D_\odot + T_\odot$ foundational interface & \textbf{Complete (grammar = HTT, 2026-03-31)} \\
Step 3 & Express $\Gamma_\text{seq}$ as type-theoretic temporal dependencies & Open \\
Step 4 & Verify Corollary 3.12 in extended system & Open \\
\bottomrule
\end{tabularx}
\end{table}

The realization of Step 2 reduces the outstanding work from 4 open steps to 3. More
importantly, the HTT foundation is no longer a specification to be built: it is a working
algebraic system (the SynthOmnicon grammar) with a proof assistant encoding
(\texttt{Synthon.lean}, \texttt{Core.lean}) and a machine-checked Frobenius theorem (Theorem 23.1,
PRIMITIVE\_THEOREMS \pilcrow{}23). Steps 3 and 4 can now be attempted within this foundation.

\begin{center}
\rule{0.5\textwidth}{0.4pt}
\end{center}

\subsection{Acknowledgments}

The author thanks the SynthOmnicon framework and its open-source development
community for providing the computational infrastructure used in this analysis.
The author would also like to thank Harry T. Larson for imparting 
the importance of catching a rising problem and never letting it go

\begin{center}
\rule{0.5\textwidth}{0.4pt}
\end{center}


\end{document}
